\chapter{Recomendaciones y Conclusiones}

\section{Conclusiones Principales}

\subsection{Análisis de Flujo de Carga}

Del análisis de flujo de carga se puede observar que ante la entrada en servicio de la generación:

\begin{itemize}
    \item La cargabilidad de todos los elementos del sistema eléctrico del área de influencia del proyecto y de alimentación (propiedad del OR) se encuentran dentro de los límites de operación.
    
    \item Las variaciones de tensión en los nodos cercanos al punto de conexión son máximas de 0.00062 p.u., lo cual está dentro de los límites permitidos.
    
    \item No existe un aumento de la cargabilidad en ninguna línea de distribución, indicando que la energía generada es consumida de manera local.
\end{itemize}

\subsection{Análisis de Cortocircuito}

El análisis de cortocircuito muestra que:

\begin{itemize}
    \item Se presenta un aumento máximo de 0.00000003 kA en la corriente de cortocircuito en el sistema, lo cual no genera la necesidad de cambios en las protecciones existentes.
    
    \item La variación máxima es menor al 0.01\% respecto al caso base.
    
    \item La tecnología de generación empleada (inversores solares) limita efectivamente la corriente de inyección hacia el sistema.
\end{itemize}

\subsection{Análisis de Protecciones}

\begin{itemize}
    \item Se validó el cumplimiento de los requisitos de protección establecidos en el Acuerdo 1322 del CNO.
    
    \item No es necesario realizar modificaciones en las protecciones existentes.
    
    \item El sistema cuenta con protección anti-isla, de sobretensión, baja tensión y de frecuencia.
\end{itemize}

\subsection{Análisis de Pérdidas}

\begin{itemize}
    \item La instalación del sistema solar fotovoltaico produce una reducción en las pérdidas técnicas del sistema.
    
    \item Se observa una reducción máxima de 0.34 kW en las pérdidas a las 12 p.m. y 3 p.m.
    
    \item Esto mejora la eficiencia global de la red de distribución.
\end{itemize}

\section{Recomendaciones}

\subsection{Medición de Calidad de Potencia}

Se recomienda que una vez entre en servicio el proyecto se realicen mediciones de calidad de potencia para determinar que el factor de potencia se encuentre dentro del rango permitido por la regulación y evitar posibles sanciones por parte del OR.

\textbf{Justificación:} Al bajar el consumo de potencia activa y permanecer igual el consumo de reactiva, el factor de potencia tiende a bajar.

\textbf{Acción correctiva:} Si una vez realizadas las mediciones se determina que el factor de potencia está fuera de los límites, se podría proceder a la instalación de un banco de condensadores que compense los reactivos necesarios.

\subsection{Uso de Resultados de Cortocircuito}

Los resultados del estudio de cortocircuito para el sistema eléctrico del generador fotovoltaico son un valioso insumo para cuando se desee realizar las siguientes actividades, teniendo en cuenta las normas ANSI e IEC:

\begin{enumerate}
    \item Evaluar y/o especificar las capacidades de interrupción de interruptores y demás equipos de seccionamiento tales como fusibles y seccionadores.
    
    \item Evaluación de saturación de transformadores de corriente.
    
    \item Diseño de mallas de puesta a tierra.
    
    \item Los niveles de cortocircuito presentados en este informe deben ser considerados para efectuar la especificación de nuevos equipos en los casos donde sean requeridos.
\end{enumerate}

\subsection{Validación en Campo}

Se recomienda realizar mediciones de campo antes y después de la puesta en servicio del proyecto con el fin de:

\begin{itemize}
    \item Verificar el cumplimiento de los límites operativos aquí establecidos.
    
    \item Validar que los niveles de tensión se mantengan dentro de los rangos establecidos.
    
    \item Confirmar que no existan sobrecargas en los elementos del sistema.
    
    \item Verificar la calidad de la potencia inyectada a la red.
\end{itemize}

\subsection{Coordinación con el Operador de Red}

Es importante mantener una coordinación permanente con ESSA para:

\begin{itemize}
    \item Garantizar que la operación del sistema se realice dentro de los límites establecidos.
    
    \item Reportar cualquier anomalía o cambio en las condiciones de operación.
    
    \item Realizar los ajustes necesarios en las protecciones si en el futuro se modifican las características del circuito.
\end{itemize}

\section{Hallazgos Positivos}

\begin{enumerate}
    \item \textbf{Bajo impacto en la red:} El proyecto tiene un impacto mínimo en el sistema de distribución existente, con variaciones de tensión y cortocircuito prácticamente imperceptibles.
    
    \item \textbf{Capacidad disponible:} El transformador de conexión opera por debajo del 60\% de su capacidad, dejando margen para futuras expansiones o incrementos de demanda.
    
    \item \textbf{Reducción de pérdidas:} La generación distribuida contribuye a la reducción de pérdidas técnicas en la red.
    
    \item \textbf{Generación local:} La totalidad de la energía generada es consumida localmente, evitando inyecciones innecesarias a la red.
    
    \item \textbf{Protecciones adecuadas:} El sistema cuenta con todas las protecciones requeridas y validadas según la normativa colombiana.
\end{enumerate}

\section{Resumen Ejecutivo}

El presente estudio de conexión simplificado para el proyecto solar fotovoltaico MAS X MENOS de 141.6 kWp en Bucaramanga, Santander, concluye que:

\begin{quote}
    \textit{``El proyecto de generación solar fotovoltaica de 141.6 kWp propuesto para el Supermercado MAS X MENOS sede Guarín puede conectarse al Sistema de Distribución Local de ESSA en el nodo 3272966 sin generar violaciones a los parámetros de operación normal de la red. La tecnología de generación empleada limita efectivamente la corriente de inyección hacia el sistema, resultando en un aumento insignificante de las corrientes de cortocircuito. Las protecciones existentes cumplen con los requisitos regulatorios y no requieren modificaciones. La instalación contribuirá a la reducción de pérdidas técnicas en la red local.''}
\end{quote}

Por lo anterior, se \textbf{APRUEBA} la conexión del proyecto propuesto, recomendando la implementación de las medidas de monitoreo y validación de calidad de potencia mencionadas en este documento.
